\chapter{Cystic Fibrosis Genetic Testing}

\section*{What Is This Test?}
Cystic-fibrosis genetic testing identifies pathogenic variants in the CFTR gene. It is used for diagnosis, carrier assessment, newborn follow-up, and selected reproductive or family evaluations.

\section*{Alternative Names}
CFTR Mutation Analysis; Cystic Fibrosis Carrier Test; CFTR Sequencing.

\section*{Why This Test Is Ordered}
Testing is performed after an abnormal newborn screen, in people with clinical features such as chronic sinopulmonary disease or pancreatic insufficiency, and for partners or relatives of a person with a known CFTR variant. It may guide eligibility for genotype-specific therapy.

\section*{How This Test Is Performed}
DNA is usually extracted from blood or a cheek sample. A common-variant panel may be followed by full CFTR sequencing and deletion/duplication analysis when the initial result is negative or incomplete.

\section*{Preparation, Risks, and Limitations}
No fasting is required. Risks are limited to sample collection. A negative limited panel does not exclude cystic fibrosis, particularly in people with ancestry underrepresented in the panel. Genetic counseling is recommended when results affect reproductive decisions.

\section*{Understanding Results}
Two disease-causing CFTR variants in trans support a molecular diagnosis, but genotype must be interpreted with sweat chloride, clinical findings, and functional testing. A single pathogenic variant usually indicates carrier status unless another variant is not detected.

\section*{References}
Cystic Fibrosis Foundation guidance on diagnosis and CFTR genetic testing.

\clearpage
\section*{Regulatory Status and Authorization Date}
Regulatory status applies to the specific assay, instrument, manufacturer, and intended use rather than to the test name alone. No single approval date can be assigned to this test category. For a commercial assay, verify the current product in the relevant regulator's database: the U.S. Food and Drug Administration (FDA) clearance, approval, or authorization; India's Central Drugs Standard Control Organisation (CDSCO) licence or permission; the European Union In Vitro Diagnostic Regulation (EU IVDR) CE certificate issued through the applicable conformity-assessment route; or the United Kingdom Medicines and Healthcare products Regulatory Agency (MHRA) registration and UKCA/CE status. Laboratory-developed tests may not have an individual FDA, CDSCO, EU IVDR, or MHRA approval date.
\clearpage
